\subsectionaddtolist{الف}
ابزارهای خودکار شناسایی میکروسرویس (مانند \lr{Mono2Micro}) معمولاً معماری‌های کاندیدا را صرفاً بر اساس تحلیل توالی دسترسی به داده‌ها، وابستگی‌های کدی و لاگ‌های اجرایی تولید می‌کنند. دلیل ناکافی بودن آن‌ها برای \lr{DDD} این است که خروجی این ابزارها فاقد درک معنایی از منطق کسب‌وکار است و خروجی آن‌ها قابلیت مدل‌سازی یا بازآرایی با مفاهیم غنی \lr{DDD} را ندارد. به عبارت دیگر، این ابزارها یک برش فنی از سیستم ارائه می‌دهند اما شکاف بزرگی بین این ساختارهای استخراج‌شده و الگوهای معنایی دامنه وجود دارد که معماران برای توسعه سیستم به آن نیاز دارند.

\vspace{0.5em}

\subsectionaddtolist{ب}

در ابزار \lr{Mono2Micro}، معیار اصلی برای کنار هم قرار دادن موجودیت‌ها و تشکیل یک \lr{Cluster}، مفهوم \lr{Transactional Similarity} است؛ به این معنا که موجودیت‌هایی که معمولاً در جریان یک تراکنش به‌صورت همزمان خوانده یا نوشته می‌شوند، با احتمال بیشتری در یک خوشه قرار می‌گیرند.

هنگامی که محققان تلاش کردند این خوشه‌ها را به مفاهیم \lr{DDD} نگاشت کنند، مشخص شد که هر خوشه به‌طور همزمان با دو مفهوم متفاوت در \lr{DDD} همپوشانی معنایی دارد. از یک سو، این خوشه‌ها از نظر رفتاری به \lr{Aggregate} شباهت دارند، زیرا در \lr{DDD} نیز \lr{Aggregate} مجموعه‌ای از اشیاء دامنه است که به‌شدت به هم وابسته‌اند و برای اعمال تغییرات، به‌صورت یک واحد تراکنشی یکپارچه عمل می‌کنند. از سوی دیگر، همین خوشه‌ها از منظر ساختاری به \lr{Bounded Context} نیز نزدیک هستند، زیرا مرزهای تعامل و وابستگی بین بخش‌های مختلف سیستم را مشخص می‌کنند و عملاً نقش جداسازی در سطح معماری و میکروسرویس را ایفا می‌نمایند.

با توجه به این ماهیت دوگانه، راهکار نهایی ارائه‌شده در مقاله این است که هر خوشه استخراج‌شده از \lr{Mono2Micro} به یک \lr{Bounded Context} مجزا تبدیل شود و درون هر \lr{Bounded Context} تنها یک \lr{Aggregate} منفرد و نسبتاً بزرگ تعریف گردد که تمامی موجودیت‌های آن خوشه را در بر بگیرد. این تصمیم به‌عنوان یک نقطه شروع برای مدل‌سازی در نظر گرفته می‌شود و نه یک طراحی نهایی ثابت، زیرا به معمار اجازه می‌دهد ابتدا ساختار کلی سیستم را بر اساس داده‌های رفتاری واقعی شکل دهد و سپس آن را به‌صورت تدریجی بهینه‌سازی کند.

مقاله همچنین توضیح می‌دهد که سایر ترکیب‌های ممکن چرا کنار گذاشته شده‌اند؛ اگر هر خوشه صرفاً به یک \lr{Aggregate} تبدیل می‌شد و همه در یک \lr{Bounded Context} بزرگ قرار می‌گرفتند، ماهیت توزیع‌شده و استقلال سرویس‌ها که هدف اصلی معماری میکروسرویس است از بین می‌رفت. از طرف دیگر، اگر هر خوشه فقط به یک \lr{Bounded Context} تبدیل می‌شد اما درون آن چندین \lr{Aggregate} مستقل تعریف می‌گردید، ارتباطات تراکنشی واقعی که باعث شکل‌گیری آن خوشه شده‌اند در مدل نهایی بازتاب دقیقی نداشتند و انسجام رفتاری سیستم از دست می‌رفت.

در ادامه، مقاله تأکید می‌کند که این نگاشت اولیه قابل تکامل است. معمار می‌تواند بر اساس دانش دامنه و تحلیل دقیق‌تر، \lr{Aggregate} اولیه را بازآرایی و به چند \lr{Aggregate} کوچک‌تر و معنایی‌تر تقسیم کند. همچنین برای مدیریت وابستگی‌ها پس از تفکیک کانتکست‌ها، به‌جای ارجاعات مستقیم بین موجودیت‌ها، از موجودیت‌های مرجع محلی استفاده می‌شود تا مرزهای \lr{Bounded Context} حفظ شود و در عین حال امکان شناسایی و بازطراحی ارتباطات بین بخش‌های مختلف سیستم برای تیم معماری فراهم بماند.

\vspace{0.5em}

\subsectionaddtolist{ج}
ابزار \lr{Context Mapper} در این پایپ‌لاین نقش موتور مدل‌سازی و تولید زبان را بر عهده دارد. این ابزار خروجی‌های \lr{Mono2Micro} را دریافت کرده و با استفاده از زبان دامنه اختصاصی خود به نام \lr{CML}، آن‌ها را به صورت خودکار به مدل‌های \lr{DDD} ترجمه می‌کند. مزیت کلیدی این رویکرد برای معمار نرم‌افزار این است که با در اختیار داشتن یک مدل واضح، متنی و بصری (\lr{CML})، می‌تواند پیچیدگی‌های سیستم استخراج‌شده را به راحتی درک کند و مرزهای کانتکست‌ها را با استفاده از الگوهای استراتژیک و تاکتیکی \lr{DDD} به شکلی چابک‌تر ویرایش، بازآرایی و بهینه‌سازی نماید.